\documentclass[10pt,a4paper]{amsart}
\usepackage[margin=19mm]{geometry}
\usepackage{amsmath,amssymb,mathtools}
\usepackage[scr=rsfs,cal=euler]{mathalfa}
\usepackage{xfrac}
\usepackage[unicode,hidelinks,bookmarks=false]{hyperref}
\newcommand{\ie}{i.e.\ }
\newcommand{\cf}{cf.\ }
\newcommand{\resp}{respectively\ }
\newcommand{\Subsection}{\subsection}
\providecommand{\nobreakdash}{\nobreak\hbox{-}}
\setlength{\emergencystretch}{2em}
\multlinegap0pt
\usepackage{orcidlink}
\DeclareMathOperator{\RE}{Re}
\newcommand{\T}{\mathbb{T}}
\def\D{B(0,1)}
\def\sub{\subseteq}
\def\LebT{\tau}
\def\lpar{\left(}
\def\rpar{\right)}
\def\into{\longrightarrow}
\def\sub{\subseteq}
\def\setdiff{\backslash}
\def\rlim{{\displaystyle\lim_{r\rightarrow 1^-}}}
\def\klim{{\displaystyle\lim_{k\rightarrow\infty}}}
\def\Bbar{\overline{B}}
 \newtheorem{theorem}{Theorem}[section]
 \newtheorem{lemma}[theorem]{Lemma}
  \newtheorem{proposition}[theorem]{Proposition}
 \newtheorem{corollary}[theorem]{Corollary}
\begin{document}
\title[Sufficient Conditions for Strong Natural Boundaries involvin]{Sufficient Conditions for Strong Natural Boundaries involving Lacunary Sequences}
\author{Rafael Reno S. Cantuba}
\date{}
\maketitle
\noindent\emph{Source and licence.} Sufficient Conditions for Strong Natural Boundaries involving Lacunary Sequences. Version arXiv:2608.17197v1 (CC BY 4.0). This material is distributed under the Creative Commons Attribution 4.0 International licence, http://creativecommons.org/licenses/by/4.0/, as recorded in the arXiv OAI licence metadata for this exact version and shown on the abstract page https://arxiv.org/abs/2608.17197v1. Full rights notice: licenses/arxiv-2608.17197-CC-BY-4.0.txt.\par\smallskip
\noindent\textit{Editorial scope.} Counted unit: the original \texttt{proposition} environment \texttt{coeffProp} at source lines 100--105 together with its complete proof at lines 107--130; the delivered document keeps source lines 96--131 so that every referenced label and every citation resolves inside it. Native editable source is the paper's own concatenated TeX; the AMS wrapper is editorial formatting only. No publisher class, upstream script or unrelated artwork is redistributed.
\section{Boundedness of coefficients}\label{CoeffSec}

For the weak regular point $1$ of $f$, the image $U(1)$ of a fixed subinterval $(a,b) \subseteq (-\delta,\delta) $ under the function $t \mapsto e^{it}$ is strictly within the arc $I_\delta$. Let $\chi_{U(1)}$ denote the characteristic function or (set) indicator function of $U(1)$. For a fixed $r \in (0,1)$ and $x \in [-\pi, \pi)$, the function $e^{i\theta} \mapsto f(re^{i(\theta+x)})$ is in the function space $C(\T)$ of all complex-valued continuous functions on $\T$, and if we denote Lebesgue measure on $\T$ by $\LebT$, then the function $e^{i\theta}\mapsto f(re^{i(\theta+x)})\chi_{U(1)}(e^{i\theta})$ is in $L^{1}(\LebT)$. By one form of the Riesz Representation Theorem, there exists a Radon measure $\iota_x$ in the Banach space $M(\T)$, the topological dual of $C(\T)$, such that $g\in C(\T)$ implies $\int_\T g\  d\iota_x= \int_{-\pi}^\pi g(e^{i\theta})f(re^{i(\theta+x)})\chi_{U(1)}(e^{i\theta})\  \frac{d\theta}{2\pi}$. The Fourier-Stieltjes transform of $\iota_x$ is the complex-valued function $\widehat{\iota_x}$ on the dual of the compact group $\T$ defined by $\widehat{\iota_x}(n_k) = \int_{\T} e^{-in_k\theta} d\iota_x(e^{i\theta})$.


\begin{proposition}\label{coeffProp}
If $1$ is a weak regular point of $f$, then for each positive integer $k$, 
\[
a_k = \frac{2\pi}{r^{n_k}(b-a)} \int_{-\pi}^\pi \widehat{\iota_x}(n_k) e^{-in_k x} \frac{dx}{2\pi}.
\]
\end{proposition}

\begin{proof}
By the definition of the measure $\iota_x$ and of Fourier-Stieltjes transform,
\begin{equation}
\widehat{\iota_x}(n_k) = \int_{a}^{b} f(re^{i(\theta+x)})e^{-in_k\theta}\ \frac{d\theta}{2\pi}.\label{FS_trans}
\end{equation}
Multiplying both sides of \eqref{FS_trans} by $e^{-in_k x}$ and integrating with respect to $x$ over $\T$,
\[
\int_{-\pi}^\pi \widehat{\iota_x}(n_k) e^{-in_k x} \frac{dx}{2\pi} = \int_{-\pi}^\pi \left( \int_{a}^{b} f(re^{i(\theta+x)}) e^{-in_k\theta} \ \frac{d\theta}{2\pi} \right) e^{-in_k x} \frac{dx}{2\pi},
\]
 to which we substitute the power series expansion for $f(re^{i(\theta+x)})$, which results to
\[
\int_{-\pi}^\pi \widehat{\iota_x}(n_k) e^{-in_k x} \frac{dx}{2\pi} = \int_{-\pi}^\pi \left( \int_{a}^{b} \sum_{m=1}^\infty a_m r^{n_m} e^{in_m(\theta+x)} e^{-in_k\theta} \ \frac{d\theta}{2\pi} \right) e^{-in_k x} \frac{dx}{2\pi}.
\]
Since $r \in (0,1)$ is fixed and $f$ is analytic on the open disk $B(0,1)$, the power series expansion $\sum_{m=1}^\infty a_m z^{n_m}$ converges absolutely and uniformly on the (compact) boundary of the disk $B(0,r)\sub B(0,1)$. Consequently, the series for $f(re^{i(\theta+x)})$ converges absolutely and uniformly with respect to $(\theta, x) \in [a,b] \times [-\pi, \pi)$. Due to this uniform convergence, term-by-term integration is justified, and Fubini's theorem may be used to obtain
\begin{align*}
\int_{-\pi}^\pi \widehat{\iota_x}(n_k) e^{-in_k x} \frac{dx}{2\pi} &= \sum_{m=1}^\infty a_m r^{n_m} \int_{-\pi}^\pi \left( \int_{a}^{b} e^{in_m\theta} e^{in_mx} e^{-in_k\theta} e^{-in_k x} \ \frac{d\theta}{2\pi} \right) \frac{dx}{2\pi}, \\
&= \sum_{m=1}^\infty a_m r^{n_m} \left( \int_{a}^{b} e^{i(n_m - n_k)\theta}\ \frac{d\theta}{2\pi} \right) \left( \int_{-\pi}^\pi e^{i(n_m - n_k)x} \frac{dx}{2\pi} \right).
\end{align*}
We now evaluate the integral over the parameter $x$. Because $(n_k)$ is a strictly increasing sequence of positive integers, the orthogonality of characters on $\T$ implies that the integral $\int_{-\pi}^\pi e^{i(n_m - n_k)x} \frac{dx}{2\pi}$ evaluates to $1$ if $m = k$, and $0$ if $m \neq k$. Thus, 
\[
\int_{-\pi}^\pi \widehat{\iota_x}(n_k) e^{-in_k x} \frac{dx}{2\pi} = a_k r^{n_k} \left( \int_{a}^{b} 1\ \frac{d\theta}{2\pi} \right) \cdot 1 = a_k r^{n_k} \frac{b-a}{2\pi}.
\]
Isolating $a_k$ yields the desired equation.
\end{proof}


\end{document}
